% 资源分配例的可行多边形：内部点(0.6,0.6)无活跃约束，全部方向可行；边上点(1.7,0.6)只有 2x1+x2<=4 活跃，可行方向为半平面 2d1+d2<=0；
% 顶点(4/3,4/3)两个约束活跃，可行方向为锥 {2d1+d2<=0, d1+2d2<=0}，由射线(1,-2)与(-2,1)张成。
\begin{tikzpicture}[x=15mm,y=15mm,
  note/.style={font=\FigureNoteFont\fontsize{7.5}{9}\selectfont,text=FigureMuted}]
  \fill[FigureBlueFill] (0,0)--(2,0)--(4/3,4/3)--(0,2)--cycle;
  \draw[FigureBlue,line width=.9pt] (0,0)--(2,0)--(4/3,4/3)--(0,2)--cycle;
  \draw[FigureYellow,line width=2pt] (2,0)--(4/3,4/3);
  \draw[->,FigureStroke,line width=.6pt] (-0.15,0)--(2.75,0) node[right]{$x_1$};
  \draw[->,FigureStroke,line width=.6pt] (0,-0.15)--(0,2.3) node[above]{$x_2$};
  % interior point
  \fill[FigureGreenFill,draw=FigureGreen,line width=.6pt] (0.6,0.6) circle (0.28);
  \fill[FigureInk] (0.6,0.6) circle (1.3pt);
  % edge point
  \fill[FigureGreenFill,draw=FigureGreen,line width=.6pt] (1.7,0.6) -- ++(116.57:0.28) arc[start angle=116.57,end angle=296.57,radius=0.28] -- cycle;
  \fill[FigureInk] (1.7,0.6) circle (1.3pt);
  % vertex
  \fill[FigureGreenFill,draw=FigureGreen,line width=.6pt] (4/3,4/3) -- ++(-63.43:0.32) arc[start angle=-63.43,end angle=-206.57,radius=0.32] -- cycle;
  \fill[FigureRed] (4/3,4/3) circle (1.5pt);
  \node[note,anchor=north] at (0.72,0.29) {内部：无活跃约束};
  \node[note,anchor=west,align=left] at (2.0,0.75) {边上：一个\\活跃约束};
  \node[note,anchor=west,align=left] at (1.45,1.62) {顶点：两个活跃约束};
  \node[note,anchor=north] at (2,-0.05) {$(2,0)$};
  \node[note,anchor=east] at (-0.03,2) {$(0,2)$};
\end{tikzpicture}
